\section{Fast Fourier Transform}

\subsection{Layerwise contraction}

We contract the FFT tensor network layerwise.
There are $2\cdot \catorder$ layers, each has a contraction demand in $\mathcal{O}(2^{\catorder})$ (exploiting diagonality of $U$).
The algorithm has a total computational demand in $\mathcal{O}(d2^{\catorder})$.

\subsection{Divide and conquer}

Called the radix-2 FFT \cite{cooley_algorithm_1965}.
Can be seen as an iterative simplification by adding a $\onesat{\catvariableof{\catorder-1}}$ and splitting it into the basis vectors.


\stantikzpic{

    \draw (-5,-3) rectangle (-3,3);
    \drawtextnode{-4}{0}{\corelabelsize $\exvector$}

    \drawhwire{-3}{2.5}{$\catvariableof{0}$}
    \drawtextnode{-2}{0.5}{\colorlabelsize $\vdots$}
    \drawhwire{-3}{-2.5}{$\catvariableof{\catorder-1}$}
    \draw (-1,-3) rectangle (1,3);
    \drawtextnode{0}{0}{\corelabelsize $F^{\catorder}$}
    \drawhwire{1}{2.5}{$\headvariableof{\catorder-1}$}
    \drawtextnode{2}{0.5}{\colorlabelsize $\vdots$}
    \drawhwire{1}{-2.5}{$\headvariableof{0}$}

    \drawtextnode{6}{0}{$=$}

    \begin{scope}[shift={(15,0)}]

        \draw (-6,-3) rectangle (-4,3);
        \drawtextnode{-5}{0}{\corelabelsize $\exvector$}

        \drawhwire{-4}{2.5}{$\catvariableof{0}$}
        \drawtextnode{-3}{1.5}{\colorlabelsize $\vdots$}
        \drawhwire{-4}{-0.5}{$\catvariableof{\catorder-2}$}

        \draw (-2,-1) rectangle (0,3);
        \drawtextnode{-1}{1}{\corelabelsize $F^{\catorder-1}$}

        \drawvariabledot{-1}{-2}
        \draw (-1,-2) -- (-1,-3);
        \drawtensorblock{-1}{-4}{2}{$\fbasis+\tbasis$}

        \draw (-4,-2) -- (5,-2);
        \drawsmalltensor{6}{-2}{$\hgate$}
        \drawhwire{7}{-2}{$\headvariableof{0}$}
        \draw (3,-2) -- (3,-1);
        \drawvariabledot{3}{-2}

        \drawhwire{0}{2.5}{$\headvariableof{\catorder-1}$}
        \drawtextnode{1}{1.5}{\colorlabelsize $\vdots$}
        \drawhwire{0}{-0.5}{$\headvariableof{1}$}
        \draw (2,-1) rectangle (4,3);
        \drawtextnode{3}{1}{\corelabelsize $U^{\catorder-1}$}
        \draw (4,2.5) -- (7,2.5);
        \drawtextnode{5}{1.5}{\colorlabelsize $\vdots$}
        \drawhwire{7}{2.5}{$\headvariableof{\catorder-1}$}
        \draw (4,-0.5) -- (7,-0.5);
        \drawhwire{7}{-0.5}{$\headvariableof{1}$}

    \end{scope}

}

Now by linearity of contractions, this splits into two terms

\stantikzpic{

    \begin{scope}[shift={(0,0)}]

        \draw (-6,-3) rectangle (-4,3);
        \drawtextnode{-5}{0}{\corelabelsize $\exvector$}

        \drawhwire{-4}{2.5}{$\catvariableof{0}$}
        \drawtextnode{-3}{1.5}{\colorlabelsize $\vdots$}
        \drawhwire{-4}{-0.5}{$\catvariableof{\catorder-2}$}

        \draw (-2,-1) rectangle (0,3);
        \drawtextnode{-1}{1}{\corelabelsize $F^{\catorder-1}$}

        \draw (-4,-2.5) -- (-3,-2.5);
        \drawsmalltensor{-2}{-2.5}{$\fbasis$}

        \drawsmalltensor{1}{-2.5}{$\fbasis$}

        \draw (2,-2.5) -- (5,-2.5);
        \drawsmalltensor{6}{-2.5}{$\hgate$}
        \drawhwire{7}{-2.5}{$\headvariableof{0}$}
        \draw (3,-2.5) -- (3,-1);
        \drawvariabledot{3}{-2.5}

        \drawhwire{0}{2.5}{$\headvariableof{\catorder-1}$}
        \drawtextnode{1}{1.5}{\colorlabelsize $\vdots$}
        \drawhwire{0}{-0.5}{$\headvariableof{1}$}
        \draw (2,-1) rectangle (4,3);
        \drawtextnode{3}{1}{\corelabelsize $U^{\catorder-1}$}
        \draw (4,2.5) -- (7,2.5);
        \drawtextnode{5}{1.5}{\colorlabelsize $\vdots$}
        \drawhwire{7}{2.5}{$\headvariableof{\catorder-1}$}
        \draw (4,-0.5) -- (7,-0.5);
        \drawhwire{7}{-0.5}{$\headvariableof{1}$}

    \end{scope}

    \drawtextnode{15}{0}{$+$}

    \begin{scope}[shift={(25,0)}]

        \draw (-6,-3) rectangle (-4,3);
        \drawtextnode{-5}{0}{\corelabelsize $\exvector$}

        \drawhwire{-4}{2.5}{$\catvariableof{0}$}
        \drawtextnode{-3}{1.5}{\colorlabelsize $\vdots$}
        \drawhwire{-4}{-0.5}{$\catvariableof{\catorder-2}$}

        \draw (-2,-1) rectangle (0,3);
        \drawtextnode{-1}{1}{\corelabelsize $F^{\catorder-1}$}

        \draw (-4,-2.5) -- (-3,-2.5);
        \drawsmalltensor{-2}{-2.5}{$\tbasis$}

        \drawsmalltensor{1}{-2.5}{$\tbasis$}

        \draw (2,-2.5) -- (5,-2.5);
        \drawsmalltensor{6}{-2.5}{$\hgate$}
        \drawhwire{7}{-2.5}{$\headvariableof{0}$}
        \draw (3,-2.5) -- (3,-1);
        \drawvariabledot{3}{-2.5}

        \drawhwire{0}{2.5}{$\headvariableof{\catorder-1}$}
        \drawtextnode{1}{1.5}{\colorlabelsize $\vdots$}
        \drawhwire{0}{-0.5}{$\headvariableof{1}$}
        \draw (2,-1) rectangle (4,3);
        \drawtextnode{3}{1}{\corelabelsize $U^{\catorder-1}$}
        \draw (4,2.5) -- (7,2.5);
        \drawtextnode{5}{1.5}{\colorlabelsize $\vdots$}
        \drawhwire{7}{2.5}{$\headvariableof{\catorder-1}$}
        \draw (4,-0.5) -- (7,-0.5);
        \drawhwire{7}{-0.5}{$\headvariableof{1}$}

    \end{scope}

}

We simplify them to


\stantikzpic{

    \draw (-6,-3) rectangle (-4,3);
    \drawtextnode{-5}{0}{\corelabelsize $\exvector$}

    \drawhwire{-4}{2.5}{$\catvariableof{0}$}
    \drawtextnode{-3}{1.5}{\colorlabelsize $\vdots$}
    \drawhwire{-4}{-0.5}{$\catvariableof{\catorder-2}$}

    \draw (-2,-1) rectangle (0,3);
    \drawtextnode{-1}{1}{\corelabelsize $F^{\catorder-1}$}

    \draw (-4,-2.5) -- (-3,-2.5);
    \drawsmalltensor{-2}{-2.5}{$\fbasis$}

    \drawsmalltensor{1}{-2.5}{$\hplusbas$}


    \drawhwire{2}{-2.5}{$\headvariableof{0}$}

    \draw (0,2.5) -- (2,2.5);
    \drawhwire{2}{2.5}{$\headvariableof{\catorder-1}$}
    \drawtextnode{1}{1.5}{\colorlabelsize $\vdots$}
    \draw (0,-0.5) -- (2,-0.5);
    \drawhwire{2}{-0.5}{$\headvariableof{1}$}


    \drawtextnode{6}{0}{$+$}

    \begin{scope}[shift={(15,0)}]

        \draw (-6,-3) rectangle (-4,3);
        \drawtextnode{-5}{0}{\corelabelsize $\exvector$}

        \drawhwire{-4}{2.5}{$\catvariableof{0}$}
        \drawtextnode{-3}{1.5}{\colorlabelsize $\vdots$}
        \drawhwire{-4}{-0.5}{$\catvariableof{\catorder-2}$}

        \draw (-2,-1) rectangle (0,3);
        \drawtextnode{-1}{1}{\corelabelsize $F^{\catorder-1}$}

        \draw (-4,-2.5) -- (-3,-2.5);
        \drawsmalltensor{-2}{-2.5}{$\tbasis$}

        \drawsmalltensor{1}{-2.5}{$\hminusbas$}
        \draw (2,-2.5) -- (4,-2.5);
        \drawhwire{4}{-2.5}{$\headvariableof{0}$}


        \drawhwire{0}{2.5}{$\headvariableof{\catorder-1}$}
        \drawtextnode{1}{1.5}{\colorlabelsize $\vdots$}
        \drawhwire{0}{-0.5}{$\headvariableof{1}$}
        \draw (2,-1) rectangle (4,3);
        \drawtextnode{3}{1}{\corelabelsize $U^{\catorder-1}$}

        \drawtextnode{5}{1.5}{\colorlabelsize $\vdots$}
        \drawhwire{4}{2.5}{$\headvariableof{\catorder-1}$}
        \drawhwire{4}{-0.5}{$\headvariableof{1}$}

    \end{scope}

}

The contraction of $\exvector$ with $\tbasisat{\catvariableof{\catorder-1}}$ respectively $\tbasisat{\catvariableof{\catorder-1}}$ are the even and the odd coordinates of $\exvector$.
Both are Fourier transformed with ${\catorder-1}$ bits, which are further divided into Fourier transforms of lower order.
In the conquer step of the FFT, the odd fourier transform gets a coordinatewise transform by $U$, which is efficient since $U$ is diagonal.
